% UET Thesis Pandoc LaTeX Template
% Customized for University of Engineering and Technology, VNU

\clearpage
\phantomsection

\setcounter{chapter}{5}
\chapter[Kết luận]{Kết luận}

\section{Kết quả đạt
được}\label{kux1ebft-quux1ea3-ux111ux1ea1t-ux111ux1b0ux1ee3c}

Khóa luận đã hoàn thành mục tiêu đề ra: thiết kế và cài đặt từ đầu (from
scratch) một hệ thống mật mã đường cong elliptic hoàn chỉnh trên trường
nguyên tố 1024-bit bằng ngôn ngữ Rust, không phụ thuộc bất kỳ thư viện
mật mã bên ngoài nào. Các đóng góp cụ thể bao gồm:

\begin{itemize}
\item
  \textbf{Lớp số học bignum an toàn bộ nhớ:} Tự định nghĩa kiểu dữ liệu
  \texttt{U1024} với phép cộng, trừ, nhân, lũy thừa và nghịch đảo. Phép
  nhân Montgomery \citep{Montgomery1985} được cài đặt trực tiếp, loại bỏ
  hoàn toàn phép chia số lớn trong vòng lặp tính toán, đồng thời bảo đảm
  an toàn bộ nhớ tuyệt đối nhờ kiến trúc Rust.
\item
  \textbf{Đường cong pairing-friendly mới 1024-bit:} Áp dụng thành công
  Phương pháp Nhân phức (CM) và thuật toán Cocks-Pinch cải tiến
  \citep{CocksPinch2001} để chủ động sinh đường cong KSS18 với bậc nhúng
  \(k = 18\), bậc nhóm \(r \approx 2^{512}\) (bảo mật 256-bit), và cấu
  trúc NTT-friendly trên cả hai trường \(\mathbb{F}_r\) và
  \(\mathbb{F}_p\). Đây là đóng góp gốc của khóa luận --- đường cong tự
  sinh, không lấy từ bất kỳ tiêu chuẩn nước ngoài nào, đồng thời kháng
  được tấn công TNFS \citep{BarbulescuDuquesne2019}, MOV
  \citep{MOV1993}, và Anomalous \citep{SmartAnomaly1999}.
\item
  \textbf{Ba sơ đồ chữ ký số chuẩn mực:} Schnorr \citep{Schnorr1989} và
  ECDSA \citep{FIPS186_4} được cài đặt đầy đủ và xác minh đúng đắn. Nền
  tảng lý thuyết của BLS \citep{BLS2001} được trình bày và tích hợp vào
  kiến trúc, sẵn sàng cho cài đặt phép ghép cặp đầy đủ trong tương lai.
\item
  \textbf{Bộ kiểm thử toàn diện:} 36/36 test PASSED, bao phủ từ số học
  cấp thấp (\texttt{U1024}, \texttt{FieldElement}) đến an toàn mật mã
  cấp cao (mô phỏng tấn công MOV, Anomalous, Invalid Curve).
\end{itemize}

\section{Đánh giá tổng quan}\label{ux111uxe1nh-giuxe1-tux1ed5ng-quan}

\subsection{So sánh với mục tiêu ban
đầu}\label{so-suxe1nh-vux1edbi-mux1ee5c-tiuxeau-ban-ux111ux1ea7u}

Nhìn lại bốn mục tiêu đặt ra tại Chương 1, khóa luận đạt được kết quả
như sau:

\begin{longtable}[]{@{}
  >{\raggedright\arraybackslash}p{(\linewidth - 2\tabcolsep) * \real{0.5263}}
  >{\raggedright\arraybackslash}p{(\linewidth - 2\tabcolsep) * \real{0.4737}}@{}}
\caption{Đối chiếu kết quả với mục tiêu đề ra}\tabularnewline
\toprule\noalign{}
\begin{minipage}[b]{\linewidth}\raggedright
Mục tiêu
\end{minipage} & \begin{minipage}[b]{\linewidth}\raggedright
Kết quả
\end{minipage} \\
\midrule\noalign{}
\endfirsthead
\toprule\noalign{}
\begin{minipage}[b]{\linewidth}\raggedright
Mục tiêu
\end{minipage} & \begin{minipage}[b]{\linewidth}\raggedright
Kết quả
\end{minipage} \\
\midrule\noalign{}
\endhead
\bottomrule\noalign{}
\endlastfoot
Tự chủ tham số đường cong & \(\checkmark\) Sinh được đường cong KSS18
1024-bit mới, không tái sử dụng hằng số bên ngoài \\
Bảo mật trước máy tính cổ điển & \(\checkmark\) Bảo mật 256-bit, vượt xa
chuẩn NIST 128-bit đến năm 2030 \\
An toàn bộ nhớ ngôn ngữ & \(\checkmark\) Rust loại bỏ toàn bộ lỗi bộ nhớ
tại thời điểm biên dịch \\
Chữ ký số vận hành được & \(\checkmark\) Schnorr và ECDSA ký/xác minh
chính xác, 36/36 test PASSED \\
\end{longtable}

\subsection{Đánh giá điểm mạnh và hạn
chế}\label{ux111uxe1nh-giuxe1-ux111iux1ec3m-mux1ea1nh-vuxe0-hux1ea1n-chux1ebf}

\textbf{Điểm mạnh:}

Hệ thống đạt được tính \emph{đúng đắn toán học} cao --- mọi công thức từ
phép cộng điểm affine, phép nhân Montgomery đến quy trình ký/xác minh
Schnorr và ECDSA đều được kiểm thử độc lập. Tính \emph{minh bạch thuật
toán} là ưu điểm nổi bật so với các thư viện ``hộp đen'' công nghiệp:
toàn bộ tham số đường cong sinh ra từ quy trình Cocks-Pinch xác định,
không có hạt giống (seed) bí ẩn, đảm bảo không thể tồn tại backdoor ẩn.

\textbf{Hạn chế:}

Hiệu năng hiện tại (\textasciitilde1.5 s/phép ký) chưa đủ cho các ứng
dụng yêu cầu độ trễ thấp như TLS hay thanh toán thời gian thực. Đây là
sự đánh đổi có chủ ý khi ưu tiên tính đúng đắn và kiến trúc phân tầng rõ
ràng trước, tối ưu hóa sau. Ngoài ra, thuật toán Double-and-Add hiện tại
tiềm ẩn nguy cơ timing side-channel do số phép \texttt{add} phụ thuộc
vào trọng số Hamming của scalar --- một điểm yếu cần giải quyết trước
khi đưa vào môi trường production.

\section{Hướng phát triển}\label{hux1b0ux1edbng-phuxe1t-triux1ec3n}

\subsection{Tối ưu hóa hiệu năng (ngắn
hạn)}\label{tux1ed1i-ux1b0u-huxf3a-hiux1ec7u-nux103ng-ngux1eafn-hux1ea1n}

\begin{enumerate}
\def\labelenumi{\arabic{enumi}.}
\item
  \textbf{Tọa độ Jacobian/Projective:} Loại bỏ phép nghịch đảo tốn kém
  (\(a^{p-2} \bmod p\)) trong mỗi bước \texttt{add}/\texttt{double} bằng
  cách biểu diễn điểm ở dạng \((X:Y:Z)\) thay vì \((x, y)\) Affine. Toàn
  bộ vòng lặp nhân vô hướng 1024 bước chỉ cần 1 phép nghịch đảo duy nhất
  ở cuối --- ước tính giảm thời gian ký từ \textasciitilde1.5 s xuống
  dưới \textbf{200 ms}.
\item
  \textbf{Sliding Window / NAF:} Giảm \textasciitilde25--30\% số phép
  \texttt{add} so với Double-and-Add cơ bản bằng cách nhóm nhiều bit
  scalar lại. Kết hợp với Jacobian, thời gian ký dự kiến đạt dưới
  \textbf{150 ms}.
\item
  \textbf{AVX-512 / SIMD assembly:} Vector hóa các phép nhân 64×64-bit
  limb trong Montgomery multiplication trên kiến trúc x86-64. Thư viện
  \texttt{blst} của Ethereum đạt tốc độ BLS12-381 signing dưới 1 ms nhờ
  kỹ thuật này --- con số tham chiếu thực tế cho hướng tối ưu tiếp theo.
\end{enumerate}

\subsection{Bảo mật cấp cao (trung
hạn)}\label{bux1ea3o-mux1eadt-cux1ea5p-cao-trung-hux1ea1n}

\begin{enumerate}
\def\labelenumi{\arabic{enumi}.}
\setcounter{enumi}{3}
\item
  \textbf{Montgomery Ladder:} Thay thế Double-and-Add bằng Montgomery
  Ladder để đạt \emph{constant-time} thực sự --- thuật toán luôn thực
  hiện đúng 1024 phép \texttt{double} và 1024 phép \texttt{add} bất kể
  giá trị scalar, loại bỏ hoàn toàn timing side-channel.
\item
  \textbf{Phép ghép cặp đầy đủ (Pairing):} Cài đặt vòng lặp Miller và
  bước lũy thừa cuối (final exponentiation) trên \(\mathbb{F}_{p^{18}}\)
  để kích hoạt BLS aggregate signatures. Đây là hướng có giá trị ứng
  dụng cao nhất --- một chữ ký BLS tổng hợp từ \(n\) người ký có thể xác
  minh bằng đúng 2 phép ghép cặp, không phụ thuộc vào \(n\).
\item
  \textbf{Multi-Scalar Multiplication (MSM --- Pippenger):} Tối ưu bước
  xác minh \([s]G + [e]Q\) bằng thuật toán Pippenger \citep{BLS12_381},
  giảm \textasciitilde50\% công việc so với hai phép nhân vô hướng tuần
  tự.
\end{enumerate}

\subsection{Mở rộng ứng dụng (dài
hạn)}\label{mux1edf-rux1ed9ng-ux1ee9ng-dux1ee5ng-duxe0i-hux1ea1n}

\begin{enumerate}
\def\labelenumi{\arabic{enumi}.}
\setcounter{enumi}{6}
\item
  \textbf{Giao thức ngưỡng (Threshold Signatures):} Dựa trên tính tổng
  hợp tự nhiên của Schnorr, cài đặt các giao thức MuSig2 hay FROST cho
  phép \(t\)-trong-\(n\) người ký tham gia, phù hợp với ứng dụng đa chữ
  ký (multi-sig) trong blockchain và hệ thống tài chính.
\item
  \textbf{Kiểm định hình thức (Formal Verification):} Sử dụng framework
  như \texttt{creusot} hoặc \texttt{kani} để chứng minh tính đúng đắn
  của các tính chất quan trọng (ví dụ: \texttt{add} thỏa mãn luật giao
  hoán và kết hợp) ở mức ngôn ngữ hình thức, nâng mức độ đảm bảo an toàn
  lên chuẩn mực cao nhất.
\end{enumerate}

Các hướng \textbf{1} và \textbf{4} là ưu tiên cao nhất vì chúng đồng
thời cải thiện cả hiệu năng lẫn bảo mật mà không đòi hỏi thay đổi kiến
trúc cấp cao, và có thể được triển khai độc lập trong phạm vi codebase
hiện tại.

\FloatBarrier
